\documentclass[11pt,a4paper]{article}
\usepackage[T1]{fontenc}
\usepackage{lmodern}
\usepackage[margin=25mm]{geometry}
\usepackage{amsmath,amssymb,amsthm,mathtools}
\usepackage{booktabs,longtable,array}
\usepackage{microtype,xcolor,enumitem,fancyhdr}
\usepackage{xurl}
\usepackage[colorlinks=true,linkcolor=blue!45!black,urlcolor=blue!55!black,citecolor=blue!45!black]{hyperref}
\usepackage{listings}
\hypersetup{
  pdftitle={Semantic Anchors for Diophantine Consistency},
  pdfauthor={AI-assisted research prepared for Vladimir Reshetnikov},
  pdfsubject={Explicit quartic Diophantine equations equivalent to the inconsistency of ZFC},
  pdfkeywords={Diophantine equations, ZFC consistency, Pell equations, C222, ProveIt}
}
\definecolor{ink}{HTML}{19364B}
\definecolor{muted}{HTML}{526575}
\pagestyle{fancy}
\fancyhf{}
\fancyhead[L]{\small\color{muted}Semantic anchors for Diophantine consistency}
\fancyhead[R]{\small\color{muted}ProveIt research}
\fancyfoot[C]{\thepage}
\setlength{\headheight}{14pt}
\setlength{\parskip}{4pt}
\setlength{\parindent}{0pt}
\setlength{\emergencystretch}{2em}
\newtheorem{theorem}{Theorem}[section]
\newtheorem{lemma}[theorem]{Lemma}
\newtheorem{proposition}[theorem]{Proposition}
\theoremstyle{remark}
\newtheorem{remark}[theorem]{Remark}
\newcommand{\N}{\mathbb N}
\newcommand{\Z}{\mathbb Z}
\newcommand{\Con}{\operatorname{Con}}
\newcommand{\ZFC}{\mathsf{ZFC}}
\newcommand{\ZF}{\mathsf{ZF}}
\newcommand{\Ccal}{\mathcal C_{222}}
\newcommand{\htop}{\operatorname{ht}}
\newcommand{\qref}[1]{\nolinkurl{#1}}
\newcommand{\one}{\mathbf 1}
\lstset{basicstyle=\ttfamily\small,breaklines=true,columns=fullflexible,frame=single,rulecolor=\color{black!15}}

\title{\color{ink}\textbf{Semantic Anchors for Diophantine Consistency}\\[5pt]
  \large An explicit quartic in 234 nonnegative integers,\\
  a 233-variable companion, and a reusable certificate reduction}
\author{Research prepared for Vladimir Reshetnikov\\
  \small AI-assisted mathematical construction and independent computational review}
\date{24 October 2026}

\begin{document}
\maketitle
\thispagestyle{empty}

\begin{abstract}
We construct two literal integer polynomials of degree four whose lack of
nonnegative integer zeros is provably equivalent in ZFC to the syntactic
consistency of ZFC. The recommended polynomial has 234 variables, 183 quadratic
residuals, 9,019 collected monomials, and coefficient height 87,741,134. A
233-variable companion has 182 residuals and 50,762 monomials. Every residual
and every coefficient is supplied, together with deterministic reconstruction
and independently checked elimination certificates.

The construction combines the repository's shared Pell bound with its
six-power modular-binomial compiler. Its additional structural idea is to
reuse an already-required, large label polynomial as the sole fully certified
Pell index. A controlled Chinese remainder lift makes that index dominate all
other required indices in the completeness proof. No additional dominance
predicate is needed in the final equation. Consecutive Pell ordinates and a
dominant digit denominator give further simplifications.

Relative to the inspected 261-variable direct construction, the recommended
equation removes 27 variables and 11 residuals, and reduces collected support
by about 71.5\%; coefficient height rises by about 2.5\%. This is a concrete
improvement of the direct numerical-calculus route. It is not a claim of
global minimality, a new variables-only universal bound, or a new complete
Lean/Rocq formalization. The proof uses the inherited numerical
\(\Ccal\)-to-refutation calibration; it neither simulates K215 nor
instantiates Jones's universal polynomial.
\end{abstract}

\vfill
\textbf{Exact comparison snapshot.}
\texttt{f4d8db99e060339f3306a04d507c35bc1d0d68c5}, ProveIt.
The main branch was changing during this investigation; all comparisons
and provenance below refer to this frozen snapshot.
\newpage
{\setlength{\parskip}{0pt}\tableofcontents}
\newpage

\section{The completed finite objects}

Throughout, \(\N=\{0,1,2,\ldots\}\). Write \(r_0,\ldots,r_{182}\) for
the \emph{literal} quadratic polynomials printed in Appendix~\ref{app:residuals}
and stored in \qref{equations/natural234/equation.json}. The recommended equation is
\begin{equation}\label{eq:main}
  \boxed{P_{234}(x_0,\ldots,x_{233})=
    \sum_{j=0}^{182}r_j(x_0,\ldots,x_{233})^2=0,
    \qquad (x_0,\ldots,x_{233})\in\N^{234}.}
\end{equation}
All indices and coefficients in this definition are fixed integers. In
particular, the subscripts do not hide an unknown horizon or a variable
number of equations. A second complete definition, with every collected
coefficient, is \qref{equations/natural234/expanded_polynomial.json}.

The companion \(P_{233}\) is the sum of the squares of the 182 residuals in
\qref{equations/natural233/equation.json}. Its full residual list is also
provided as a plain-text equation and a separate LaTeX input. Neither equation
contains a variable exponent, a factorial operator, a binomial operator, a
bounded quantifier, or an unspecified universal-polynomial constant.

\begin{theorem}\label{thm:main}
For the two fixed polynomials just specified, ordinary mathematics formalizable
in ZFC proves
\[
 \ZFC\vdash
 \bigl(\neg\exists\mathbf x\in\N^{234}\;P_{234}(\mathbf x)=0\bigr)
 \longleftrightarrow \Con(\ZFC),
\]
and the same statement with \(P_{233}\) and \(\N^{233}\).
Here \(\Con(\ZFC)\) means that the fixed effective proof calculus for ZFC
has no finite proof of contradiction.
\end{theorem}

The proof establishes the equivalent positive assertion: a zero exists if and
only if a finite ZFC refutation exists. It uses no assumption of consistency.
It does not decide either polynomial's solvability. The full end-to-end
argument is an ordinary proof, with precisely identified inherited logical
interfaces; the accompanying programs verify the finite polynomial objects
and transformations, not a newly generated internal-ZFC proof term.

\subsection{A measured comparison}

\begin{center}\small
\begin{tabular}{lrrrrr}
\toprule
Natural-number quartic & Variables & Rows & Degree & Terms & Height\\
\midrule
Shared-binomial parent &319&254&4&10,283&86,842,335\\
Shared Pell bound &261&194&4&31,679&85,566,955\\
Its sparse companion &292&225&4&6,140&83,933,334\\
\textbf{Recommended here} &\textbf{234}&\textbf{183}&4&\textbf{9,019}&87,741,134\\
Smaller-variable companion &\textbf{233}&182&4&50,762&88,674,690\\
\bottomrule
\end{tabular}
\end{center}
``Rows'' counts quadratic residuals before the outer sum of squares; ``terms''
and ``height'' refer to the \emph{collected single polynomial}. These metrics
are checked from the actual JSON polynomials, including coordinate liveness.
The two prior arrivals are cited and pinned in the provenance record
\cite{sharedbinomial,cap261}. The sparse 292-variable point remains a
smaller-support tradeoff. The recommended new equation has one more variable
than its companion and approximately 5.6 times fewer monomials.

The repository already has constructions with many fewer variables at other
degrees, and separate universal-polynomial-based quartics with fewer variables
than these direct equations. In the literature, nine natural unknowns are
known to suffice for general Diophantine representations, and Sun's September
2026 preprint gives the corresponding ten-integer-unknown result
\cite{sun2026}. These are different optimization objectives. No lower bound
proving that 233, 234, or any listed degree is optimal is supplied here.

\section{The idea: an existing witness can certify the others}

Three pieces of earlier work fit together:
\begin{enumerate}[leftmargin=*]
\item A single indexed Pell solution can bound many unindexed Pell solutions.
One congruence for the sum of squared index errors then checks every index
at once \cite{cap261}.
\item Many large binomial numerators are needed only modulo a fixed modulus.
The modular compiler obtains them from Pell residues and retains just six
exact arithmetic powers \cite{sharedbinomial,residuecompiler}.
\item The source contains an already-required polynomial \(J\), the sum of its
possible rule labels. A Chinese remainder representative may be increased by
the entire CRT period without changing any of the encoded local rows.
\end{enumerate}

The first two pieces cannot be combined by deleting their separate checks
without changing the common size budget. The modular compiler introduced
additional Pell indices that were absent from its exact-power budget. All
those indices must enter the new budget. Also, the Pell growth certificate
must constrain a lower bound \(L\), while the actual base \(A\ge L\) remains
free to meet the common-root congruence.

The further reduction is to make \(J\) itself the one certified cap index.
In completeness, one encoded input is increased by \(N^2\). Because \(J\)
contains a quadratic constructor of that input, it then exceeds every other
required index. The aggregate range budget changes from \(9N^2\) to
\(10N^2\). In soundness, no such lift is assumed: the cap and pooled error
equation themselves force the other indices to be correct. Consequently no
additional inequality asserting \(J\)'s dominance appears in the polynomial.

This use of a semantic witness as a certificate for the other witnesses is
the principal additional idea of this iteration. The construction and its
measured improvement are new relative to the inspected packages; no claim
of literature-wide priority is made.

\section{Pell arithmetic and the simultaneous index theorem}

\subsection{Elementary Pell facts}

For an integer \(A\ge2\), define
\[
 X_0(A)=1,\quad Y_0(A)=0,\qquad
 X_{n+1}=AX_n+(A^2-1)Y_n,\quad Y_{n+1}=X_n+AY_n.
\]
The identity
\[
 X_n(A)^2-(A^2-1)Y_n(A)^2=1
\]
holds by induction. Every nonnegative integer solution of this Pell equation
is a unique pair \((X_n(A),Y_n(A))\). One direct proof applies the inverse
map
\[
 (X,Y)\longmapsto (AX-(A^2-1)Y,\ AY-X).
\]
For \(Y>1\), both coordinates are nonnegative and the second strictly
decreases: \((A-1)Y<X<AY\). The case \(Y=1\) is \((A,1)\), which
descends to \((1,0)\). This proves the classification by induction.

The same recurrence gives, for the applicable positive indices,
\begin{align}
 X_n(A)&\equiv1\pmod{A-1}, &Y_n(A)&\equiv n\pmod{A-1},\label{eq:pellresidue}\\
 X_n(A)&>(A-1)Y_n(A), &Y_{n+1}(A)&>(2A-1)Y_n(A),\label{eq:pellgrowth}
\end{align}
and \(Y_n(A)\ge n\). These are elementary finite inductions, available
inside ZFC without a consistency hypothesis.

\subsection{The inherited one-index certificate}

We retain one complete instance of the positive indexed-Pell criterion.
For a requested index \(H\), a pair \((X_*,Y_*)\), and natural auxiliary
variables, it consists of
\begin{align}
 X_*^2-(A^2-1)Y_*^2&=1,\nonumber\\
 U^2-(A^2-1)V^2&=1,\nonumber\\
 S^2-(C^2-1)T^2&=1,\nonumber\\
 C&=1+4Y_*(d+1), &V&=Y_*^2(e+1),\label{eq:indexcriterion}\\
 C&=A+Uf, &S&=X_*+Ug,\nonumber\\
 T&=H+4Y_*h.\nonumber
\end{align}
Together with \(H\le Y_*\), the inherited criterion identifies
\((X_*,Y_*)=(X_H(A),Y_H(A))\). Its completeness statement supplies the
Pell pairs, divisibility, and modular equalities for canonical pairs.
This is an imported arithmetic
theorem, not an inference from a finite search. Its exact source interface
is \qref{DiophantineConsistency.PellRepresentation.criterion_iff_pell};
the retained statements, proofs, and source pins are included in the package.
The aggregate row below supplies \(H\le Y_*\), so a separate slack for
that hypothesis is unnecessary.

The displayed rows use \emph{nonnegative} congruence quotients, which need
one additional check beyond the modular statement of the imported theorem.
In this application \(H=J\ge2\), so \(Y_*\ge Y_2(A)=2A\).
The imported auxiliary base has \(C>1\) and
\(C\equiv1\pmod{4Y_*}\), hence \(C\ge1+4Y_*>A\).
Also \(T\equiv H\pmod{4Y_*}\), with \(0<H\le Y_*<4Y_*\), gives
\(T\ge H\ge2\). Pell classification at base \(C\) then yields
\(T\ge2C\), and
\[
 S>(C-1)T\ge4Y_*\cdot2C>AY_*>X_*.
\]
Thus the quotients in \(C-A=Uf\), \(S-X_*=Ug\), and
\(T-H=4Y_*h\) are nonnegative. The positive auxiliary ordinate and its
divisibility by \(Y_*^2\) give \(V=Y_*^2(e+1)\). This justifies the
oriented polynomial certificate in full.

\subsection{A cap with a separate size budget}

\begin{lemma}[Generalized simultaneous cap]\label{lem:cap}
Let \(W\ge3\), \(2\le H<W\), \(A>W+1\), and \(A-1>W^3\). Let
\(r\le W-2\) and \(0\le k_i<W\). Suppose \(Y_*=Y_H(A)\), and each
\((x_i,y_i)\) is a nonnegative solution of
\(x_i^2-(A^2-1)y_i^2=1\). If natural \(s,q\) satisfy
\begin{equation}\label{eq:pool}
 Y_*=H+1+s+\sum_{i=1}^r y_i,
 \qquad
 \sum_{i=1}^r(y_i-k_i)^2=(A-1)q,
\end{equation}
then \((x_i,y_i)=(X_{k_i}(A),Y_{k_i}(A))\) for every \(i\).
Conversely, if all \(k_i<H\), the canonical pairs extend to a natural
solution of \eqref{eq:pool}.
\end{lemma}

\begin{proof}
Pell classification writes \(y_i=Y_{j_i}(A)\). The first equation gives
\(y_i<Y_H(A)\), hence \(0\le j_i<H<W\). By \eqref{eq:pellresidue},
the second equation implies that \(A-1\) divides
\[
 E=\sum_i(j_i-k_i)^2.
\]
But
\[
 0\le E\le r(W-1)^2\le(W-2)(W-1)^2<W^3<A-1.
\]
Therefore \(E=0\), so every \(j_i=k_i\). Notice that the proof did not
assume \(k_i<H\); it obtains that fact as a consequence of the equations.

For completeness take the canonical pairs. Each \(y_i-k_i\) is a
nonnegative multiple of \(A-1\), so the displayed quotient \(q\) is
natural. Since \(k_i<H\),
\[
 Y_H(A)-\sum_iY_{k_i}(A)
 >(2A-1-r)Y_{H-1}(A)>W>H.
\]
Thus \(s=Y_H(A)-H-1-\sum_iY_{k_i}(A)\) is natural.
\end{proof}

This version separates the actual cap index \(H\) from the size budget
\(W\). That separation is exactly what permits the semantic anchor
\(H=J\).

\subsection{Consecutive coordinates need no added orientation guard}

The fourteen ordinary seeds use natural coordinates \((y_i,z_i)\) and
a product wire \(u_i\), with
\begin{equation}\label{eq:consecutive}
 u_i=y_iz_i,\qquad y_i^2+z_i^2-2Au_i=1.
\end{equation}
Their existing power or modular-output row has the form
\begin{equation}\label{eq:orient}
 n_i y_i=z_i+R_i,
 \qquad R_i\ge0,\quad 1\le n_i<A.
\end{equation}
If \(y_i=0\), \eqref{eq:consecutive} gives \(z_i=1\), contradicting
\eqref{eq:orient}. Therefore \(y_i>0\), and
\(z_i\le n_iy_i<Ay_i\). Set \(x_i=Ay_i-z_i>0\). Direct expansion gives
\[
 x_i^2-(A^2-1)y_i^2=y_i^2+z_i^2-2Ay_iz_i=1.
\]
Thus the cap lemma applies. After it recovers the index, the recurrence
identifies \(z_i=Y_{k_i-1}(A)\). Conversely, for \(k_i\ge1\), these
consecutive ordinates satisfy all the displayed identities and
\[
 X_{k_i}(A)+(n_i-A)Y_{k_i}(A)=n_iy_i-z_i.
\]
The sole full anchor pair retains the ordinary \((X_*,Y_*)\) coordinates.
Every seed's output incidence and the bound \(n_i<A\) are checked in the
semantic audit; an unconnected conic would not suffice.

\section{Recovering all powers before using their meanings}

\subsection{The literal shared growth frame}

Let \(\mathcal A\) contain the distinct base, exponent, and output atoms of
the six exact-power requests, together with every modular upper-index atom.
It includes \(J\) and \(q\). The graph sets
\begin{equation}\label{eq:frame}
 W=16+\sum_{a\in\mathcal A}a,
 \qquad Z=z+1,
 \qquad
 L^2-\big((W+1)^2-1\big)(WZ)^2=1,
 \qquad A=L+\sigma.
\end{equation}
All atoms and witnesses are natural. In particular, every intended index,
power input, and proposed power output is less than \(W\). There are
fourteen ordinary seeds, so \(14\le W-2\). Also \(q+1<W\), since
\(q\in\mathcal A\) and the offset is 16; a separate \(q+1\) atom is
unnecessary.

Pell classification of the third equation writes
\(L=X_e(W+1)\), \(WZ=Y_e(W+1)\). Modulo \(W\), the recurrence gives
\(Y_e(W+1)\equiv e\), so \(W\mid e\). Since \(WZ>0\), \(e\ge W\).
Consequently
\begin{equation}\label{eq:largeA}
 A\ge L\ge X_W(W+1)\ge(W+1)^W,
 \qquad A-1>W^3.
\end{equation}
This is established from polynomial rows before interpreting any proposed
arithmetic power or factorial.

For completeness choose \(e=W\), then \(Z=Y_W(W+1)/W\ge1\).
The value \(L\) is fixed, while \(A\) can be any integer at least \(L\).
That freedom is essential for the modular-root congruence in the next section.

\subsection{The six exact requests}

Using outer natural variables \(t_0,x,y\) and further semantic wires, put
\begin{equation}\label{eq:semantic}
 t=t_0+16,\quad R=2+t+x+y+F,\quad K=2R,\quad B=b_0+2.
\end{equation}
The six exact-power requests are
\begin{equation}\label{eq:sixpowers}
 V_0=(F^2+1)^t,\quad N=B^K,\quad q=N^K,\quad
 D_t=q^t,\quad D_y=q^{y+1},\quad D_x=q^{x+1}.
\end{equation}
These expressions describe meanings proved for ordinary polynomial graphs.
They are not exponentiation primitives in the delivered equation.
The additional denominator \(D_{t+1}=qD_t\) uses one multiplication.

For a request \(m=n^k\), the graph defines a natural modulus
\[
 T_n=2An-n^2-1
\]
and imposes the residue relation
\[
 X_k(A)+(n-A)Y_k(A)=m+T_nh.
\]
With consecutive coordinates the left side is \(ny_k-z_k\).
The principal indices are certified first by Lemma~\ref{lem:cap}.
The modulus row itself excludes \(n=0\); all exponents are positive by
their literal definitions. From \eqref{eq:largeA} and the frame bounds,
\[
 T_n>W^W+W,\qquad 0\le m<W,\qquad 0<n^k<W^W.
\]
For example, \(T_n\ge2A-W^2-1\), and
\(W^W>W^2+W+1\) for \(W\ge3\), which proves the strict modulus bound.
Substitution \(s=n-A\) in \((A+s)^k=X_k+Y_ks\), modulo \(T_n\), gives
\(X_k+(n-A)Y_k\equiv n^k\pmod{T_n}\). The two small residues must be
equal. This proves all six powers simultaneously.

For completeness the same expression is positive, by
\(X_k>(A-1)Y_k\) and \(n\ge1\). Since its residue is the true output
\(m<T_n\), its quotient \(h\) is nonnegative.

\section{One root congruence for all eleven modular requests}

\subsection{The dominant denominator is available early}

The source contains the natural ceiling rows
\begin{equation}\label{eq:ceiling}
 V_0+B=BF+r+1,\qquad B=r+s+1,
\end{equation}
and endpoint rows
\begin{equation}\label{eq:endpoints}
 x=F+1+(y+1)u_0,\quad y=F+1+r_0,\quad
 x=1+\bigl(1+(t+1)y\bigr)u_t.
\end{equation}
After \eqref{eq:sixpowers} is decoded, \(V_0\ge1\). The ceiling rows give
\(BF=V_0+s\ge1\), hence \(F\ge1\). Thus \(x,y\ge2\), and the last
endpoint forces \(u_t\ge1\). In particular,
\begin{equation}\label{eq:dominates}
 x\ge2+(t+1)y>\max(t+1,y).
\end{equation}
This argument does not yet use \(F=t!\).

Set
\begin{equation}\label{eq:rootmod}
 M=ND_x=Nq^{x+1},\qquad b=q+1,
 \qquad 2Ab=b^2+1+Mh_0.
\end{equation}
Every required range modulus \(Nq^k\), for
\(k\in\{t,t+1,y+1,x+1\}\), divides \(M\). The exact-binomial modulus
\(q^{t+1}=qD_t\) also divides \(M\). This dominant-denominator
observation is inherited from the residue-compiler arrival
\cite{residuecompiler} and is used here in the merged-binomial architecture.

\subsection{Exact low residues from Pell pairs}

The last equation in \eqref{eq:rootmod} implies
\((b-A)^2\equiv A^2-1\pmod M\). Therefore for every \(n\ge0\),
\begin{equation}\label{eq:modpell}
 \rho_n=X_n(A)+(b-A)Y_n(A)\equiv b^n\pmod M.
\end{equation}
Moreover \(\rho_n>0\): it is 1 at \(n=0\), and exceeds
\((b-1)Y_n(A)>0\) at positive indices.

For an exact binomial \(C=\binom nk\), with \(d=q^k\), use
\begin{equation}\label{eq:exactdigit}
 \rho_n=r+dC+qd h,\qquad d=r+s+1,\qquad q=C+o+1.
\end{equation}
For a divisibility request \(N\mid\binom nk\), use
\begin{equation}\label{eq:divdigit}
 \rho_n=r+Nd h,\qquad d=r+s+1.
\end{equation}
If all coefficients \(\binom nj\), \(0\le j\le k\), are smaller than
\(q\), the binomial theorem and \eqref{eq:modpell} prove that these rows
have exactly the stated meanings. For the exact request, the residue
\(r+dC\) lies in \([0,qd)\); for divisibility, the residue is less than
\(d\). Positivity of \(\rho_n\) makes the reconstructed quotients natural.

There is always a suitable base \(A\). Since \(N\mid q\),
\(\gcd(b,M)=1\). Thus \(\gcd(2b,M)\) is 1 or 2. In the latter case
\(b\) is odd and \(b^2+1\) is even. The linear congruence
\(2bA\equiv b^2+1\pmod M\) is soluble, and its solution class contains
arbitrarily large positive integers. Choose one above \(L\), large enough
that \(h_0\ge0\). All semantic data have already been chosen at this point.

\section{The factorial and CRT binomial are the same object}

\subsection{The initial digit bound is independent of factorial semantics}

The one exact binomial is intended to be
\[
 B=\binom{F^2+t}{t}.
\]
Before making that interpretation, \(B\ge2\) is already a literal shift.
Hence the six decoded powers give
\[
 N\ge2^K,\qquad q=B^{K^2}\ge2^{4R^2}.
\]
For \(n=F^2+t\) and \(0\le j\le t\),
\[
 \binom nj\le n^t\le(R^2+R)^t\le2^{2Rt}<2^{2R^2}<q.
\]
Here \(R^2+R\le4^R\), proved by a short induction, and \(t<R\).
Thus the exact digit is carry-free and identifies \(B\) before any use
of its factorial or CRT interpretation.

\begin{lemma}[Squared-radius factorial]\label{lem:factorial}
For \(t\ge6\) and a positive integer \(F\),
\[
 F=\left\lceil\frac{(F^2+1)^t}{\binom{F^2+t}{t}}\right\rceil
 \quad\Longleftrightarrow\quad F=t!.
\]
\end{lemma}
\begin{proof}
Write
\[
 H_t(F)=t!\prod_{i=1}^t\frac{F^2+1}{F^2+i}.
\]
The asserted ceiling gives \(H_t(F)\le F\le t!\). If \(F<t\),
monotonicity of each factor in \(F^2\ge1\) gives
\(H_t(F)\ge2^t/(t+1)>t>F\), a contradiction.
For \(F\ge t\), the product inequality
\(\prod_i(1-a_i)\ge1-\sum_i a_i\), with
\(a_i=(i-1)/(F^2+i)\), gives \(H_t(F)>t!/2\).
If \(F<t^2\), then \(t!/2>t^2>F\), again a contradiction for \(t\ge6\).
Consequently \(F\ge t^2\), and
\[
 H_t(F)>t!\left(1-\frac1{2F}\right).
\]
If \(t!\ge F+1\), the right side is at least
\((F+1)(1-1/(2F))>F\), which is impossible. Thus \(F=t!\).
Conversely, for \(F=t!\ge t^2\), the same inequality gives
\(F-1/2<H_t(F)\le F\), so its ceiling is \(F\).
\end{proof}

The literal rows \eqref{eq:ceiling} assert precisely this ceiling once
\(V_0=(F^2+1)^t\) and the exact binomial have been recovered. The lemma
is credited to earlier bootstrap and shared-binomial work
\cite{sharedbinomial}; it is included to make the bootstrap order explicit.

Set \(T=F^2\). For \(1\le i\le t\), the numbers
\[
 d_i=1+\frac{T}{i}
\]
are integers, and
\begin{equation}\label{eq:crtproduct}
 \prod_{i=1}^t d_i=\binom{F^2+t}{t}=B.
\end{equation}
Each \(d_i\equiv1\pmod F\), so every prime divisor exceeds \(t\).
If a prime divided both \(d_i,d_j\), it would divide \(i-j\); therefore
the \(d_i\) are pairwise coprime. This is why the same exact binomial
serves the factorial identity and the CRT construction.

\section{Reusing the label polynomial as the cap}\label{sec:anchor}

\subsection{The inherited local numerical calculus}

The source calculus has truth label 0 and contradiction label 1, with Cantor
constructors
\[
 I(a,b)=2(a+b)(a+b+1)+4b,\qquad
 M(a,b)=I(a,b)+1,\qquad A_0(a,b)=I(a,b)+2,
\]
and the remaining fixed constructors and nine axiom rows supplied in the
retained literal source \qref{direct_calculus.py}. Define \(J\) to be the
sum of the nine axiom-conclusion polynomials, \(a,b,I(a,b),A_0(a,b)\), and
1. Thus, for every natural input,
\begin{equation}\label{eq:Jgrowth}
 J(a,b,c,d)\ge I(a,b)\ge2a^2.
\end{equation}

The eight local columns are \((a,b,c,d,S,S',u,v)\), bounded respectively
by \((t,t,t,t,y,y,x,x)\). Their nine independent encoded representatives
are \(A_a,A_b,A_c,A_d,A_S,A_{S'},A_u,A_v,C\). The last is the local
row index. The remaining upper binomial indices are \(F^2+t\) and
\(J(A_a,A_b,A_c,A_d)\). We use
\[
 H=J(A_a,A_b,A_c,A_d)
\]
as the sole full Pell anchor. The fourteen other distinct indices are
\[
 t,\ K,\ y+1,\ x+1,\ F^2+t,
 \ A_a,A_b,A_c,A_d,A_S,A_{S'},A_u,A_v,C.
\]
There is no extra ordinary seed at \(H\); the modular request for \(J\)
uses the full anchor pair directly.
In soundness, the encoded constructor inputs \(A_a,A_b\) are literal
positive shifts, so \(H\ge I(1,1)=16\); this supplies the anchor's
positive-index hypotheses before any CRT interpretation.

\subsection{The CRT lift that makes the cap available}

The inherited completeness argument chooses each positive representative
at most \(N^2\), satisfying the desired local residues modulo
\(d_i^{2K}\). By \eqref{eq:crtproduct}, their common period is
\[
 \prod_i d_i^{2K}=B^{2K}=N^2.
\]
Replace only
\begin{equation}\label{eq:lift}
 A_a\longmapsto A_a+N^2.
\end{equation}
Every local polynomial congruence, including those for \(J\) and the rule
evaluation, is preserved modulo each \(d_i^{2K}\). The shifted coordinate
lies in \((N^2,2N^2]\), and the other eight stay at most \(N^2\). Hence
the literal aggregate range bound can be
\begin{equation}\label{eq:tenbudget}
 C+\sum_{j\in\{a,b,c,d,S,S',u,v\}}A_j+g=10N^2.
\end{equation}
It still concerns nine representatives; 10 is the bound coefficient,
not the number of representatives.

The five non-CRT indices listed above are at most \(N\):
\(F^2+t<R^2+R\le4^R\le N\), \(K=2R<4^R\le N\), and the remaining
three are bounded by \(R+1\le N\). The nine CRT indices are at most
\(2N^2\). Meanwhile, \eqref{eq:Jgrowth} gives
\[
 H>2N^4>2N^2.
\]
Thus every ordinary seed index is strictly smaller than \(H\), exactly
the condition needed for the completeness direction of Lemma~\ref{lem:cap}.
The semantic data, including \(H\), are chosen before \(W,L,A\).

\subsection{Why no extra dominance condition is required}

For soundness the raw graph gives \(H<W\) and \(k_i<W\) through the
frame definition. The aggregate cap gives \(j_i<H\) for each decoded
Pell index. The squared-error bound is consequently valid even without
assuming \(k_i<H\). It then proves \(j_i=k_i\), which establishes
\(k_i<H\) as a consequence. Completeness uses the lift
\eqref{eq:lift} to produce suitable witnesses. This is a change in the
choice of existential witnesses, not a claimed bijection preserving every
old witness coordinate.

\subsection{All remaining digits stay carry-free}

For an ordinary representative \(n\le10N^2\) and
\(1\le k\le R-1\),
\[
 \binom nj\le n^k\le10^kN^{2k}
 <2^{4R}N^{2k}\le N^{2k+2}\le N^K=q
 \qquad(0\le j\le k).
\]
The strict inequality uses \(10<16\), and \(N\ge2^{2R}\).
The supplied source-majorant certificate gives
\[
 J(A_a,A_b,A_c,A_d)\le(n+1)^\ell,
 \qquad \ell=2^{22},
\]
when the four inputs are at most \(n\). Since \(N\ge11\),
\(10N^2+1\le N^3\). Also, at \(t=16\),
\(3\ell(t+1)\le2t!\), and this inequality persists by induction.
Therefore
\[
 \binom{J(A)}j\le J(A)^{t+1}
 \le N^{3\ell(t+1)}<N^K=q,
 \qquad 0\le j\le t+1.
\]
The new aggregate coefficient therefore does not invalidate any digit
interpretation. The derived value \(J(A)\) is not incorrectly inserted
into the \(10N^2\) aggregate bound.

\section{From arithmetic zeros to finite refutations}

This section states the inherited arithmetic-to-calculus interface explicitly
and explains its proof. Its detailed source and proofs are retained in
\qref{upstream/shared_binomial/direct/frozen_round14}.

\subsection{Prime-power range extraction}

For any prime \(p\), \(a\ge0\), and integers \(k,e\ge1\), the elementary implication
\[
 p^e\mid\binom ak\quad\Longrightarrow\quad a\bmod p^e<k
\]
is used. If \(a\ge k\), Legendre's valuation formula expresses the
valuation of the binomial as a sum of zero-or-one terms. If \(a\bmod p^e\ge k\), the interval
\([a-k+1,a]\) contains no multiple of \(p^e\), every term at levels
\(e,e+1,\ldots\) vanishes, leaving valuation at most \(e-1\).
This proves the implication; the case \(a<k\) is immediate.
Its converse is false in general and is not used.

For each CRT row choose a prime divisor \(p\) of \(d_i\), set
\(m=v_p(d_i)\), and use precision \(p^{Km}\). The index congruence
\(B\mid C+F^2+1\) and range condition \(N\mid\binom Ct\) identify the
small index residue as \(i-1\). The eight range conditions identify
bounded local values. A tenth range condition,
\(N\mid\binom{J(A)}{t+1}\), supplies the label bound.

The literal positive arithmetic-circuit majorant bounds the local rule
polynomial and its label sum by \(R^E\), with \(E=2^{23}\). Since
\(R\ge16!>2^{44}\), elementary growth gives \(R^E<2^R<2^K\).
The nonnegative local rule value is divisible by \(p^{Km}\ge2^K\), so it
is zero. The derived label range is similarly an actual bound, not merely
a congruence. Thus every bounded local source row is recovered.

In the reverse direction one chooses representatives modulo
\(d_i^{2K}\), rather than merely \(d_i^K\). For a range binomial, one
falling-factorial numerator factor has valuation at least \(2Km\),
while division by \(k!\) loses at most \(k-1\le K-1\). The remaining
valuation is at least \(Km\). This establishes the needed divisibilities,
including the derived polynomial \(J\). The lift \eqref{eq:lift}
preserves that stronger congruence precision.

\subsection{A finite Horn-flow certificate}

Let \(w(j)=1+jF\), where \(F=t!\). The positive weights for
\(1\le j\le t\) are pairwise coprime: a prime dividing two would
divide their index difference, hence would be at most \(t\), while
each weight is 1 modulo every prime divisor of \(F\).
The local source branches express
\[
 S'\,w(\text{conclusion})
   =S\,w(\text{premise}_1)w(\text{premise}_2),
\]
with truth 0 assigned unit weight and with the literal guarded-axiom
condition when required. The two beta remainder equations identify the
next \(S\) with the preceding \(S'\). Endpoints give initial accumulator
\(F+1=w(1)\) and final accumulator 1. Multiplying the local equations gives
\[
 \prod w(\text{conclusions})
  =w(1)\prod w(\text{premises}).
\]
Prime valuations recover the exact multiplicity balance for every positive
label. If 1 were outside the closure of the chosen rules, every rule
concluding outside the closure would have a premise outside it. Counting
outside conclusions and premises would give a nonpositive difference,
contradicting the required difference 1. Cycles in the selected rule
multiset do not invalidate this argument.

Conversely, expand a finite proof DAG into a finite proof tree and choose
\(t\ge16\) above its size, parameters, and finitely many label sums.
Pad by truth rows. The pending proof multiset gives the accumulator
weights. Choose \(y\) above their values and with the necessary factorial
divisibility; CRT supplies \(x\) for the beta remainders. This produces
the bounded local source data. All later much larger arithmetic witnesses
are chosen after this finite horizon, so the proof does not circularly
bound its horizon by its own accumulator size.

\subsection{The logical endpoint and the exact proof order}

The inherited numerical-calculus endpoint is
\qref{KEndpoint.c222_iff_zf_refutation}. Its statement relates
\(\Ccal\)'s contradiction label directly to a finite syntactic ZF
refutation. The second endpoint,
\qref{ConstructibleChoice.zf_refutation_iff_zfc_refutation}, supplies
the ZF/ZFC consistency calibration. The numerical endpoint has no
machine-bootstrap or tape-coverage hypothesis. No K215 transition table
is used by this constructor.

The full soundness order is:
\begin{enumerate}[leftmargin=*]
\item Establish the natural arithmetic-gate meanings, all literal shifts,
and the enlarged frame \eqref{eq:frame}.
\item Recover the exponential lower bound on \(A\), the anchor's exact
index \(J\), and the fourteen seed indices by the pooled error equation.
For consecutive seeds, their output rows first establish orientation.
\item Recover the six exact powers. Obtain \(F\ge1\) from the ceiling
rows, and then the dominant-denominator inequalities.
\item Decode the exact factorial binomial with its independent carry bound;
the fixed-point lemma then gives \(F=t!\).
\item Decode all remaining binomial ranges and apply prime-power CRT to
recover every bounded local source row.
\item Apply Horn-flow soundness and the two inherited logical endpoints.
\end{enumerate}
For completeness, start with the refutation, choose the bounded source data,
then \(F,B,R,K,N,q\) and the CRT representatives. Perform the one-coordinate
lift, compute \(J\), all other indices, and \(W\). Choose \(L\) and a
congruence-compatible \(A\ge L\). Supply the canonical Pell sequences,
the anchor witnesses, cap slacks, and exact digit quotients. Every remaining
gate is evaluated in dependency order. The checked polynomial projection
then produces a zero of either final quartic. This proves
Theorem~\ref{thm:main}.

\section{Literal compilation, verification, and limitations}

\subsection{Two independent finite transformations}

The semantic graph is a finite list of quadratic polynomial rows and
explicit arithmetic gates. A checked affine pass replaces small nonnegative
affine gates by their formulas and removes unreachable definitional gates.
The common normalized source has 344 variables and 293 rows. The independent
export checker reconstructs the entire pass without importing the producer.

A second certificate removes a variable only when an equation is monic in
that variable and its restoration is an integer polynomial proven
nonnegative on the natural orthant. Most restorations have nonnegative
coefficients. Two require an explicit rational sum-of-squares certificate.
The checker expands each square exactly over \(\mathbb Q\), verifies the
identity and positivity of its coefficients, substitutes the restoration
into every occurrence, and checks the final variable order and residuals.
Reversing the substitutions preserves the natural domain; substitution
preserves every equation. All residual degrees stay at most two.

For \(P_{234}\), 110 such eliminations are made while retaining the anchor
ordinate. For \(P_{233}\), one more affine elimination removes that
ordinate and makes the expanded polynomial substantially denser. The
decision to retain it in the recommended equation is a measured size
tradeoff, not an extra logical assumption.

\subsection{What the verification establishes}

The release supplies independent checks of:
\begin{itemize}[leftmargin=*]
\item the exact changed semantic rows: split scale, one anchor, fourteen
seed conics, all output incidences, pooled congruence, aggregate bound,
six exact powers, eleven modular digit requests, and dominant modulus;
\item the inherited source prefix and its permitted changes, against the
frozen shared-binomial parent;
\item affine restoration, dependency pruning, every nonnegative contraction,
every collected coefficient, exact degree, and use of every variable;
\item deliberately altered data, including coefficients, gate maps,
coordinates, missing frame atoms, and wrong index/output bindings;
\item reproducibility from the included, pinned source, and the full
LaTeX transcription of the principal residuals.
\end{itemize}
The checks use explicit exceptions and are also run under Python's
optimization mode. Separate bounded arithmetic controls exercise the
Pell identities, factorial fixed point, index-error separation, and digit
extraction. These finite controls supplement the unbounded proofs; they
are not substitutes for them.

The supplied logical-interface files record dependencies and exact statement
scope. They are not a self-contained rebuild of the entire Lean development.
This release does not claim a fresh end-to-end kernel proof, a sealed
object-language ZFC proof, a worldwide novelty result, or an optimality
theorem. In particular, proving the displayed equivalence does not prove
\(\Con(\ZFC)\).

\subsection{An unrestricted-integer equation}

For integer \(a,b,c\), put
\[
 g(a,b,c)=a^2+b^2+c(c+1).
\]
It is always nonnegative and represents every natural integer. Indeed,
Legendre's three-square theorem applied to \(4n+1\) gives one odd square
and two even squares, hence \(n=a^2+b^2+c(c+1)\).
Therefore
\[
 P_{234}\bigl(g(a_0,b_0,c_0),\ldots,g(a_{233},b_{233},c_{233})\bigr)=0
\]
is an explicit degree-eight polynomial in 702 unrestricted integer
unknowns with the same consistency equivalence. The analogous companion
uses 699 integers. The literal substitution recipes and exact evaluator
are supplied. These are domain conversions; they do not improve the
repository's specialized unrestricted-integer constructions.

\section{Research controls and the next useful target}

\subsection{A cap bound cannot replace the exponential growth frame}

It would be tempting to retain only \(A-1>W^3\), because that inequality
suffices to identify all Pell indices. It does not suffice to identify
power outputs. Take
\[
 n=2,\quad k=24,\quad m=1,\quad W=29,\quad A=2^{22}+1.
\]
Then \(A-1>W^3\), but
\[
 T_n=2An-n^2-1=2^{24}-1,\qquad n^k\equiv1\pmod{T_n}.
\]
The true index-24 Pell pair therefore passes the congruence with the false
power output \(m=1\); the quotient is nonnegative. The much stronger
growth frame is doing essential work in the present construction.

\subsection{A mathematically valid reduction can worsen the final quartic}

We tested replacing \(q=N^K\) by a cofinal Pell lower-bound certificate.
With \(r=NK\), impose
\[
 q=r(z+1),\qquad P^2-\big((r+1)^2-1\big)q^2=1.
\]
Pell classification and reduction modulo \(r\) give an index at least
\(r\), hence \(N\mid q\) and \(q\ge N^K\). These are enough for the
digit arguments. Nevertheless, on a matched standard-coordinate source,
removing one exact-power request reduced the raw variable count by one
and \emph{increased} the projected count from 234 to 236: three other
eliminations were lost. The selected equation retains the cheaper exact
power. The matched experiment and its proof are recorded in the package.

The useful next target is the repeated arithmetic cost after this compiler
has reduced the number of fully indexed pairs to one. Another generic proof
checker, or a new name for a known universal representation, is unlikely to
improve this concrete endpoint without a measured certificate reduction.
The present package isolates the changed rows and semantic interfaces so
that further experiments can be compared without re-proving the source
calculus on every iteration.

\section{Reproduction and file map}

Python 3.10 or later and the standard library suffice for the mathematical
data checks. From the package root:
\begin{lstlisting}
python3 -B code/verify_release.py
python3 -B -O code/verify_release.py
python3 -B code/verify_release.py --rebuild
\end{lstlisting}
The last command reconstructs the selected equations in a temporary
directory and compares the exact generated bytes. It uses only the source
included in this archive. No proof search or solution search is performed.

Compile the report with a standard TeX Live installation:
\begin{lstlisting}
pdflatex -interaction=nonstopmode -halt-on-error report.tex
pdflatex -interaction=nonstopmode -halt-on-error report.tex
\end{lstlisting}

\begin{longtable}{@{}p{.38\textwidth}p{.58\textwidth}@{}}
\toprule
File or directory & Purpose\\\midrule\endhead
\qref{report.tex}, \qref{report.pdf} & This report and the complete principal residual appendix.\\
\qref{equations/natural234/} & Recommended polynomial, expansion, projection certificate, and metrics.\\
\qref{equations/natural233/} & Smaller-variable companion and the corresponding exact data.\\
\qref{equations/source/} & Common semantic graph and normalized 344-variable source.\\
\qref{equations/integer/} & Literal 702- and 699-integer substitution recipes.\\
\qref{code/} & Constructor, independent checkers, evaluator, and transcription generator.\\
\qref{proofs/} & Independent arithmetic reviews and inherited proof statements.\\
\qref{upstream/} & Frozen compiler and source-calculus dependency closure, retaining licenses.\\
\qref{provenance/} & Snapshot, original archive hashes, copied-file mappings, and logical interfaces.\\
\qref{verification/} & Actual run receipts and measured comparison data.\\
\qref{SHA256SUMS} & Hashes of every delivered file except this checksum list itself.\\
\bottomrule
\end{longtable}

\begin{thebibliography}{9}
\bibitem{repo} V. Reshetnikov, \emph{ProveIt}, repository snapshot
\qref{f4d8db99e060339f3306a04d507c35bc1d0d68c5}, 24 October 2026.
\url{https://github.com/VladimirReshetnikov/ProveIt/tree/f4d8db99e060339f3306a04d507c35bc1d0d68c5/Computability/DiophantineConsistency}.
\bibitem{cap261} ProveIt incoming research package,
\emph{One Pell bound and one error congruence},
\qref{ConZFC_Shared_Pell_Bound_261_2026-10-24.zip}, in the pinned
repository's \qref{docs/incoming}. Original bytes and imported interfaces
are authenticated in this release's provenance.
\bibitem{sharedbinomial} ProveIt incoming research package,
\emph{Shared binomial certificates},
\qref{ConZFC_Shared_Binomial_Certificates_2026-10-24.zip}, especially
\qref{direct/PROOF.md} and the complete \qref{equation319} data.
The included \qref{frozen_round14} directory retains the earlier factorial,
CRT, Horn-flow, and arithmetic-compilation sources.
\bibitem{residuecompiler} ProveIt incoming research package,
\emph{Residue compiler},
\qref{ConZFC_Residue_Compiler_2026-10-24.zip}, especially the
\qref{DominantModularBuilder} in \qref{code/modular_factorial.py}.
\bibitem{sun2026} Z.-W. Sun, \emph{Ten unknowns for Hilbert's tenth problem
over the integers}, arXiv:2609.01594v2, 2 September 2026.
\url{https://arxiv.org/abs/2609.01594}.
\bibitem{bayerdavid} J. Bayer and M. David,
\emph{A Formal Proof of Complexity Bounds on Diophantine Equations},
ITP 2025, LIPIcs 352, 3:1--3:18.
\url{https://doi.org/10.4230/LIPIcs.ITP.2025.3}.
\bibitem{fgo} A. Formisano, I. Gavazzi, and E. G. Omodeo,
\emph{Diophantine Modeling of Provability in Algebraic Logic}, CILC 2025.
\url{https://ceur-ws.org/Vol-4003/paper13.pdf}.
Their relation-algebra route is relevant prior art, rather than an
unexplored alternative introduced here.
\end{thebibliography}

\clearpage
\appendix
\section{All 183 literal residuals of the recommended equation}\label{app:residuals}

The following is an automatically generated, complete transcription of
\qref{equations/natural234/equation.json}. Every variable is a nonnegative
integer. Each \(r_j\) is a polynomial of total degree at most two; the
equation is \(\sum_{j=0}^{182}r_j^2=0\). Named source labels are omitted
from this typeset list to keep the formula readable; they remain in the
authoritative JSON and plain-text version. Line breaks have no mathematical
meaning. No residual, coefficient, or monomial is suppressed.

% Complete generated transcription. Do not edit by hand.
\begingroup\small
\allowdisplaybreaks[4]
\begin{align*}
r_{0}&=-x_{2}^{2}+x_{3}
\end{align*}
\begin{align*}
r_{1}&=1-2x_{2}-x_{2}x_{4}+x_{4}+x_{5}-x_{6}
\end{align*}
\begin{align*}
r_{2}&=1+x_{4}-x_{6}-x_{7}
\end{align*}
\begin{align*}
r_{3}&=-17-x_{0}-x_{0}x_{2}-x_{0}x_{9}-17x_{2}-17x_{9}+x_{11}
\end{align*}
\begin{align*}
r_{4}&=-1+x_{1}-x_{2}-x_{2}x_{8}-2x_{8}-x_{8}x_{9}
\end{align*}
\begin{align*}
r_{5}&=-1+x_{1}-x_{10}-x_{10}x_{11}
\end{align*}
\begin{align*}
r_{6}&=2+x_{3}-x_{4}x_{22}+x_{20}-2x_{22}
\end{align*}
\begin{align*}
r_{7}&=-6-5x_{12}-x_{12}^{2}-2x_{12}x_{13}-5x_{13}-x_{13}^{2}+x_{23}
\end{align*}
\begin{align*}
r_{8}&=-72-68x_{13}-16x_{13}^{2}-16x_{13}x_{23}-34x_{23}-4x_{23}^{2}+x_{24}
\end{align*}
\begin{align*}
r_{9}&=-30-11x_{12}-x_{12}^{2}-4x_{12}x_{24}-22x_{24}-4x_{24}^{2}+x_{25}
\end{align*}
\begin{align*}
r_{10}&=-380-312x_{24}-64x_{24}^{2}-32x_{24}x_{25}-78x_{25}-4x_{25}^{2}+x_{26}
\end{align*}
\begin{align*}
r_{11}&=-6-5x_{13}-x_{13}^{2}-2x_{13}x_{14}-5x_{14}-x_{14}^{2}+x_{27}
\end{align*}
\begin{align*}
r_{12}&=-6-5x_{12}-x_{12}^{2}-2x_{12}x_{14}-5x_{14}-x_{14}^{2}+x_{28}
\end{align*}
\begin{align*}
r_{13}&=-414-22x_{12}-2x_{12}^{2}-16x_{12}x_{14}-8x_{12}x_{27}-200x_{13}-32x_{13}^{2}\\
&\quad -64x_{13}x_{14}-32x_{13}x_{23}-32x_{13}x_{28}-320x_{14}-64x_{14}^{2}-32x_{14}x_{23}\\
&\quad -32x_{14}x_{27}-32x_{14}x_{28}-100x_{23}-8x_{23}^{2}-16x_{23}x_{28}-52x_{27}-8x_{27}^{2}\\
&\quad -108x_{28}-8x_{28}^{2}+x_{29}
\end{align*}
\begin{align*}
r_{14}&=-1344-800x_{13}-128x_{13}^{2}-256x_{13}x_{14}-128x_{13}x_{23}-128x_{13}x_{28}\\
&\quad -864x_{14}-128x_{14}^{2}-128x_{14}x_{23}-128x_{14}x_{28}-400x_{23}-32x_{23}^{2}\\
&\quad -64x_{23}x_{28}-432x_{28}-32x_{28}^{2}-2x_{29}-2x_{29}^{2}+x_{30}
\end{align*}
\begin{align*}
r_{15}&=-247-168x_{12}-32x_{12}^{2}-64x_{12}x_{13}-32x_{12}x_{27}-32x_{12}x_{28}-188x_{13}\\
&\quad -32x_{13}^{2}-32x_{13}x_{27}-32x_{13}x_{28}-2x_{23}-92x_{27}-8x_{27}^{2}-16x_{27}x_{28}\\
&\quad -84x_{28}-8x_{28}^{2}+x_{31}
\end{align*}
\begin{align*}
r_{16}&=-960-672x_{12}-128x_{12}^{2}-256x_{12}x_{13}-128x_{12}x_{27}-128x_{12}x_{28}\\
&\quad -736x_{13}-128x_{13}^{2}-128x_{13}x_{27}-128x_{13}x_{28}-368x_{27}-32x_{27}^{2}\\
&\quad -64x_{27}x_{28}-336x_{28}-32x_{28}^{2}-2x_{31}-2x_{31}^{2}+x_{32}
\end{align*}
\begin{align*}
r_{17}&=-276-30x_{12}-2x_{12}^{2}-16x_{12}x_{14}-8x_{12}x_{27}-30x_{13}-2x_{13}^{2}\\
&\quad -16x_{13}x_{14}-8x_{13}x_{28}-272x_{14}-64x_{14}^{2}-32x_{14}x_{27}-32x_{14}x_{28}\\
&\quad -68x_{27}-8x_{27}^{2}-68x_{28}-8x_{28}^{2}+x_{33}
\end{align*}
\begin{align*}
r_{18}&=-552-120x_{13}-8x_{13}^{2}-64x_{13}x_{14}-32x_{13}x_{28}-544x_{14}-128x_{14}^{2}\\
&\quad -128x_{14}x_{28}-272x_{28}-32x_{28}^{2}-2x_{33}-2x_{33}^{2}+x_{34}
\end{align*}
\begin{align*}
r_{19}&=-247-4x_{12}-352x_{13}-128x_{13}^{2}-64x_{13}x_{23}-64x_{13}x_{27}-92x_{23}-8x_{23}^{2}\\
&\quad -16x_{23}x_{27}-84x_{27}-8x_{27}^{2}-2x_{28}+x_{35}
\end{align*}
\begin{align*}
r_{20}&=-960-1408x_{13}-512x_{13}^{2}-256x_{13}x_{23}-256x_{13}x_{27}-368x_{23}-32x_{23}^{2}\\
&\quad -64x_{23}x_{27}-336x_{27}-32x_{27}^{2}-2x_{35}-2x_{35}^{2}+x_{36}
\end{align*}
\begin{align*}
r_{21}&=-56-60x_{13}-16x_{13}^{2}-16x_{13}x_{23}-30x_{23}-4x_{23}^{2}+x_{37}
\end{align*}
\begin{align*}
r_{22}&=-30-11x_{12}-x_{12}^{2}-4x_{12}x_{37}-22x_{37}-4x_{37}^{2}+x_{38}
\end{align*}
\begin{align*}
r_{23}&=-380-312x_{37}-64x_{37}^{2}-32x_{37}x_{38}-78x_{38}-4x_{38}^{2}+x_{39}
\end{align*}
\begin{align*}
r_{24}&=-30-11x_{13}-x_{13}^{2}-4x_{13}x_{39}-22x_{39}-4x_{39}^{2}+x_{40}
\end{align*}
\begin{align*}
r_{25}&=-349-420x_{13}-128x_{13}^{2}-128x_{13}x_{23}-208x_{23}-32x_{23}^{2}-2x_{27}+x_{41}
\end{align*}
\begin{align*}
r_{26}&=-1368-1664x_{13}-512x_{13}^{2}-512x_{13}x_{23}-832x_{23}-128x_{23}^{2}-2x_{41}-2x_{41}^{2}\\
&\quad +x_{42}
\end{align*}
\begin{align*}
r_{27}&=-56-60x_{12}-16x_{12}^{2}-16x_{12}x_{28}-8x_{12}x_{42}-30x_{28}-4x_{28}^{2}-4x_{28}x_{42}\\
&\quad -15x_{42}-x_{42}^{2}+x_{43}
\end{align*}
\begin{align*}
r_{28}&=-6-5x_{14}-x_{14}^{2}-2x_{14}x_{15}-5x_{15}-x_{15}^{2}+x_{44}
\end{align*}
\begin{align*}
r_{29}&=-224-26x_{12}-2x_{12}^{2}-16x_{12}x_{15}-8x_{12}x_{44}-288x_{14}-128x_{14}^{2}\\
&\quad -64x_{14}x_{27}-64x_{14}x_{28}-88x_{15}-32x_{15}^{2}-32x_{15}x_{44}-76x_{27}-8x_{27}^{2}\\
&\quad -16x_{27}x_{28}-68x_{28}-8x_{28}^{2}-44x_{44}-8x_{44}^{2}+x_{45}
\end{align*}
\begin{align*}
r_{30}&=-640-1152x_{14}-512x_{14}^{2}-256x_{14}x_{27}-256x_{14}x_{28}-304x_{27}-32x_{27}^{2}\\
&\quad -64x_{27}x_{28}-272x_{28}-32x_{28}^{2}-2x_{45}-2x_{45}^{2}+x_{46}
\end{align*}
\begin{align*}
r_{31}&=-241-168x_{12}-32x_{12}^{2}-64x_{12}x_{13}-32x_{12}x_{27}-32x_{12}x_{28}-184x_{13}\\
&\quad -32x_{13}^{2}-32x_{13}x_{27}-32x_{13}x_{28}-x_{14}-92x_{27}-8x_{27}^{2}-16x_{27}x_{28}\\
&\quad -84x_{28}-8x_{28}^{2}+x_{47}
\end{align*}
\begin{align*}
r_{32}&=-42-52x_{14}-16x_{14}^{2}-16x_{14}x_{27}-26x_{27}-4x_{27}^{2}+x_{48}
\end{align*}
\begin{align*}
r_{33}&=-30-11x_{14}-x_{14}^{2}-4x_{14}x_{48}-22x_{48}-4x_{48}^{2}+x_{49}
\end{align*}
\begin{align*}
r_{34}&=-306-35x_{13}-x_{13}^{2}-16x_{13}x_{48}-4x_{13}x_{49}-280x_{48}-64x_{48}^{2}\\
&\quad -32x_{48}x_{49}-70x_{49}-4x_{49}^{2}+x_{50}
\end{align*}
\begin{align*}
r_{35}&=-4556-4320x_{48}-1024x_{48}^{2}-512x_{48}x_{49}-128x_{48}x_{50}-1080x_{49}-64x_{49}^{2}\\
&\quad -32x_{49}x_{50}-270x_{50}-4x_{50}^{2}+x_{51}
\end{align*}
\begin{align*}
r_{36}&=-30-11x_{14}-x_{14}^{2}-4x_{14}x_{51}-22x_{51}-4x_{51}^{2}+x_{52}
\end{align*}
\begin{align*}
r_{37}&=-306-35x_{13}-x_{13}^{2}-16x_{13}x_{51}-4x_{13}x_{52}-280x_{51}-64x_{51}^{2}\\
&\quad -32x_{51}x_{52}-70x_{52}-4x_{52}^{2}+x_{53}
\end{align*}
\begin{align*}
r_{38}&=-1028-672x_{12}-128x_{12}^{2}-256x_{12}x_{13}-128x_{12}x_{27}-128x_{12}x_{28}\\
&\quad -736x_{13}-128x_{13}^{2}-128x_{13}x_{27}-128x_{13}x_{28}-368x_{27}-32x_{27}^{2}\\
&\quad -64x_{27}x_{28}-336x_{28}-32x_{28}^{2}-2x_{47}-2x_{47}^{2}-32x_{51}-8x_{52}-2x_{53}\\
&\quad +x_{54}
\end{align*}
\begin{align*}
r_{39}&=-264-128x_{51}-32x_{52}-8x_{53}-2x_{54}-2x_{54}^{2}+x_{55}
\end{align*}
\begin{align*}
r_{40}&=-2-3x_{12}-x_{12}^{2}-2x_{12}x_{55}-3x_{55}-x_{55}^{2}+x_{56}
\end{align*}
\begin{align*}
r_{41}&=-12-28x_{55}-16x_{55}^{2}-16x_{55}x_{56}-14x_{56}-4x_{56}^{2}+x_{57}
\end{align*}
\begin{align*}
r_{42}&=-30-11x_{14}-x_{14}^{2}-4x_{14}x_{57}-22x_{57}-4x_{57}^{2}+x_{58}
\end{align*}
\begin{align*}
r_{43}&=-306-35x_{12}-x_{12}^{2}-16x_{12}x_{57}-4x_{12}x_{58}-280x_{57}-64x_{57}^{2}\\
&\quad -32x_{57}x_{58}-70x_{58}-4x_{58}^{2}+x_{59}
\end{align*}
\begin{align*}
r_{44}&=-4556-4320x_{57}-1024x_{57}^{2}-512x_{57}x_{58}-128x_{57}x_{59}-1080x_{58}-64x_{58}^{2}\\
&\quad -32x_{58}x_{59}-270x_{59}-4x_{59}^{2}+x_{60}
\end{align*}
\begin{align*}
r_{45}&=-90-76x_{14}-16x_{14}^{2}-16x_{14}x_{27}-16x_{14}x_{60}-38x_{27}-4x_{27}^{2}\\
&\quad -8x_{27}x_{60}-38x_{60}-4x_{60}^{2}+x_{61}
\end{align*}
\begin{align*}
r_{46}&=-4556-4320x_{51}-1024x_{51}^{2}-512x_{51}x_{52}-128x_{51}x_{53}-1080x_{52}-64x_{52}^{2}\\
&\quad -32x_{52}x_{53}-270x_{53}-4x_{53}^{2}+x_{62}
\end{align*}
\begin{align*}
r_{47}&=-682-70x_{13}-2x_{13}^{2}-32x_{13}x_{60}-8x_{13}x_{61}-592x_{60}-128x_{60}^{2}\\
&\quad -64x_{60}x_{61}-148x_{61}-8x_{61}^{2}-2x_{62}+x_{63}
\end{align*}
\begin{align*}
r_{48}&=-x_{63}-x_{63}^{2}+x_{64}
\end{align*}
\begin{align*}
r_{49}&=-42-52x_{13}-16x_{13}^{2}-8x_{13}x_{14}-16x_{13}x_{23}-13x_{14}-x_{14}^{2}-4x_{14}x_{23}\\
&\quad -26x_{23}-4x_{23}^{2}+x_{65}
\end{align*}
\begin{align*}
r_{50}&=-2278-2264x_{13}-800x_{13}^{2}-800x_{13}x_{23}-160x_{13}x_{65}-70x_{14}-2x_{14}^{2}\\
&\quad -32x_{14}x_{62}-8x_{14}x_{64}-1132x_{23}-200x_{23}^{2}-80x_{23}x_{65}-592x_{62}\\
&\quad -128x_{62}^{2}-64x_{62}x_{64}-148x_{64}-8x_{64}^{2}-228x_{65}-8x_{65}^{2}+x_{66}
\end{align*}
\begin{align*}
r_{51}&=-2712-280x_{14}-8x_{14}^{2}-128x_{14}x_{62}-32x_{14}x_{64}-2368x_{62}-512x_{62}^{2}\\
&\quad -256x_{62}x_{64}-592x_{64}-32x_{64}^{2}-2x_{66}-2x_{66}^{2}+x_{67}
\end{align*}
\begin{align*}
r_{52}&=-2-3x_{67}-x_{67}^{2}+x_{68}
\end{align*}
\begin{align*}
r_{53}&=-x_{2}-x_{2}x_{17}+x_{69}
\end{align*}
\begin{align*}
r_{54}&=-x_{16}^{2}+2x_{16}x_{17}-x_{17}^{2}+x_{70}
\end{align*}
\begin{align*}
r_{55}&=-2x_{26}x_{69}-4x_{69}+x_{71}
\end{align*}
\begin{align*}
r_{56}&=-x_{16}^{2}+2x_{16}x_{17}+2x_{16}x_{71}-x_{17}^{2}-2x_{17}x_{71}-x_{71}^{2}+x_{72}
\end{align*}
\begin{align*}
r_{57}&=-x_{30}x_{69}+x_{73}
\end{align*}
\begin{align*}
r_{58}&=-x_{16}^{2}+2x_{16}x_{17}+2x_{16}x_{73}-x_{17}^{2}-2x_{17}x_{73}-x_{73}^{2}+x_{74}
\end{align*}
\begin{align*}
r_{59}&=-x_{32}x_{69}+x_{75}
\end{align*}
\begin{align*}
r_{60}&=-x_{16}^{2}+2x_{16}x_{17}+2x_{16}x_{75}-x_{17}^{2}-2x_{17}x_{75}-x_{75}^{2}+x_{76}
\end{align*}
\begin{align*}
r_{61}&=-x_{34}x_{69}+x_{77}
\end{align*}
\begin{align*}
r_{62}&=-x_{16}^{2}+2x_{16}x_{17}+2x_{16}x_{77}-x_{17}^{2}-2x_{17}x_{77}-x_{77}^{2}+x_{78}
\end{align*}
\begin{align*}
r_{63}&=-x_{36}x_{69}+x_{79}
\end{align*}
\begin{align*}
r_{64}&=-x_{16}^{2}+2x_{16}x_{17}+2x_{16}x_{79}-x_{17}^{2}-2x_{17}x_{79}-x_{79}^{2}+x_{80}
\end{align*}
\begin{align*}
r_{65}&=-4x_{13}x_{69}-2x_{40}x_{69}-4x_{69}+x_{81}
\end{align*}
\begin{align*}
r_{66}&=-x_{16}^{2}+2x_{16}x_{17}+2x_{16}x_{81}-x_{17}^{2}-2x_{17}x_{81}-x_{81}^{2}+x_{82}
\end{align*}
\begin{align*}
r_{67}&=-4x_{42}x_{69}-2x_{43}x_{69}+x_{83}
\end{align*}
\begin{align*}
r_{68}&=-x_{16}^{2}+2x_{16}x_{17}+2x_{16}x_{83}-x_{17}^{2}-2x_{17}x_{83}-x_{83}^{2}+x_{84}
\end{align*}
\begin{align*}
r_{69}&=-x_{46}x_{69}+x_{85}
\end{align*}
\begin{align*}
r_{70}&=-x_{16}^{2}+2x_{16}x_{17}+2x_{16}x_{85}-x_{17}^{2}-2x_{17}x_{85}-x_{85}^{2}+x_{86}
\end{align*}
\begin{align*}
r_{71}&=-x_{12}^{2}+2x_{12}x_{14}-x_{14}^{2}+x_{87}
\end{align*}
\begin{align*}
r_{72}&=-x_{13}^{2}+2x_{13}x_{14}-x_{14}^{2}+x_{88}
\end{align*}
\begin{align*}
r_{73}&=-x_{87}x_{88}+x_{89}
\end{align*}
\begin{align*}
r_{74}&=-4-4x_{15}-x_{15}^{2}+x_{90}
\end{align*}
\begin{align*}
r_{75}&=-x_{89}^{2}+2x_{89}x_{90}-x_{90}^{2}+x_{91}
\end{align*}
\begin{align*}
r_{76}&=-2x_{68}x_{69}-4x_{69}+x_{92}
\end{align*}
\begin{align*}
r_{77}&=-x_{16}^{2}+2x_{16}x_{17}+2x_{16}x_{92}-x_{17}^{2}-2x_{17}x_{92}-x_{92}^{2}+x_{93}
\end{align*}
\begin{align*}
r_{78}&=-x_{2}-x_{2}x_{13}+x_{94}
\end{align*}
\begin{align*}
r_{79}&=-1-x_{17}-x_{17}x_{94}-x_{94}+x_{95}
\end{align*}
\begin{align*}
r_{80}&=-x_{2}-x_{2}x_{12}+x_{96}
\end{align*}
\begin{align*}
r_{81}&=-4x_{2}-4x_{2}x_{13}-2x_{2}x_{23}+x_{97}
\end{align*}
\begin{align*}
r_{82}&=-1-x_{16}-x_{16}x_{96}-x_{96}+x_{98}
\end{align*}
\begin{align*}
r_{83}&=-x_{97}x_{98}-x_{98}+x_{99}
\end{align*}
\begin{align*}
r_{84}&=-x_{95}^{2}+2x_{95}x_{99}-x_{99}^{2}+x_{100}
\end{align*}
\begin{align*}
r_{85}&=-6x_{2}-4x_{2}x_{13}-2x_{2}x_{23}+x_{101}
\end{align*}
\begin{align*}
r_{86}&=-1-x_{17}-x_{17}x_{101}-x_{101}+x_{102}
\end{align*}
\begin{align*}
r_{87}&=-1-x_{16}-x_{16}x_{94}-x_{94}+x_{103}
\end{align*}
\begin{align*}
r_{88}&=-x_{102}^{2}+2x_{102}x_{103}-x_{103}^{2}+x_{104}
\end{align*}
\begin{align*}
r_{89}&=-x_{70}x_{72}+x_{105}
\end{align*}
\begin{align*}
r_{90}&=-x_{74}x_{105}+x_{106}
\end{align*}
\begin{align*}
r_{91}&=-x_{76}x_{106}+x_{107}
\end{align*}
\begin{align*}
r_{92}&=-x_{78}x_{107}+x_{108}
\end{align*}
\begin{align*}
r_{93}&=-x_{80}x_{108}+x_{109}
\end{align*}
\begin{align*}
r_{94}&=-x_{82}x_{109}+x_{110}
\end{align*}
\begin{align*}
r_{95}&=-x_{84}x_{110}+x_{111}
\end{align*}
\begin{align*}
r_{96}&=-x_{86}x_{111}+x_{112}
\end{align*}
\begin{align*}
r_{97}&=-x_{91}x_{112}-x_{93}x_{112}+x_{113}
\end{align*}
\begin{align*}
r_{98}&=-x_{100}x_{113}+x_{114}
\end{align*}
\begin{align*}
r_{99}&=-2-2x_{2}-x_{2}x_{20}-2x_{9}-x_{9}x_{20}-x_{20}+x_{115}
\end{align*}
\begin{align*}
r_{100}&=-1-x_{18}-x_{18}x_{115}-x_{115}+x_{116}
\end{align*}
\begin{align*}
r_{101}&=-2-x_{2}-x_{2}x_{19}-x_{9}-x_{9}x_{19}-2x_{19}-x_{19}x_{115}-x_{115}\\
&\quad +x_{117}
\end{align*}
\begin{align*}
r_{102}&=2-4x_{1}+2x_{1}^{2}-2x_{1}x_{16}-2x_{1}x_{17}-2x_{1}x_{116}-2x_{1}x_{117}+2x_{16}\\
&\quad +x_{16}^{2}+2x_{16}x_{116}+2x_{17}+x_{17}^{2}+2x_{17}x_{117}-x_{21}x_{118}\\
&\quad +x_{104}x_{114}+2x_{116}+x_{116}^{2}+2x_{117}+x_{117}^{2}
\end{align*}
\begin{align*}
r_{103}&=-9-x_{12}-x_{13}-x_{14}-x_{15}-x_{16}-x_{17}-x_{18}-x_{19}-x_{20}\\
&\quad +10x_{21}^{2}-x_{119}
\end{align*}
\begin{align*}
r_{104}&=-2-x_{4}-x_{4}x_{169}-x_{4}x_{170}-3x_{169}-x_{169}x_{171}-3x_{170}\\
&\quad -x_{170}x_{171}-x_{171}+x_{176}+x_{177}
\end{align*}
\begin{align*}
r_{105}&=-132-4x_{0}-4x_{0}x_{122}-3x_{1}-3x_{1}x_{122}-5x_{2}-5x_{2}x_{122}-2x_{3}\\
&\quad -2x_{3}x_{122}-2x_{4}-2x_{4}x_{122}-x_{5}-x_{5}x_{122}-3x_{9}-3x_{9}x_{122}\\
&\quad -2x_{12}-2x_{12}x_{122}-14x_{13}-14x_{13}x_{122}-x_{14}-x_{14}x_{122}-x_{15}\\
&\quad -x_{15}x_{122}-x_{16}-x_{16}x_{122}-x_{17}-x_{17}x_{122}-x_{18}\\
&\quad -x_{18}x_{122}-x_{19}-x_{19}x_{122}-x_{20}-x_{20}x_{122}-x_{21}\\
&\quad -x_{21}x_{122}-4x_{23}-4x_{23}x_{122}-2x_{26}-2x_{26}x_{122}-x_{30}\\
&\quad -x_{30}x_{122}-x_{32}-x_{32}x_{122}-x_{34}-x_{34}x_{122}-x_{36}\\
&\quad -x_{36}x_{122}-2x_{40}-2x_{40}x_{122}-4x_{42}-4x_{42}x_{122}-2x_{43}\\
&\quad -2x_{43}x_{122}-x_{46}-x_{46}x_{122}-2x_{68}-2x_{68}x_{122}-132x_{122}\\
&\quad -x_{122}x_{169}-x_{122}x_{170}-x_{122}x_{171}-x_{122}x_{201}-x_{122}x_{202}\\
&\quad -x_{122}x_{214}-x_{122}x_{215}+x_{123}-x_{169}-x_{170}-x_{171}-x_{201}\\
&\quad -x_{202}-x_{214}-x_{215}
\end{align*}
\begin{align*}
r_{106}&=-4x_{0}x_{123}-3x_{1}x_{123}-5x_{2}x_{123}-2x_{3}x_{123}-2x_{4}x_{123}-x_{5}x_{123}\\
&\quad -3x_{9}x_{123}-2x_{12}x_{123}-14x_{13}x_{123}-x_{14}x_{123}-x_{15}x_{123}\\
&\quad -x_{16}x_{123}-x_{17}x_{123}-x_{18}x_{123}-x_{19}x_{123}-x_{20}x_{123}\\
&\quad -x_{21}x_{123}-4x_{23}x_{123}-2x_{26}x_{123}-x_{30}x_{123}-x_{32}x_{123}\\
&\quad -x_{34}x_{123}-x_{36}x_{123}-2x_{40}x_{123}-4x_{42}x_{123}-2x_{43}x_{123}\\
&\quad -x_{46}x_{123}-2x_{68}x_{123}-133x_{123}-x_{123}x_{169}-x_{123}x_{170}\\
&\quad -x_{123}x_{171}-x_{123}x_{201}-x_{123}x_{202}-x_{123}x_{214}-x_{123}x_{215}\\
&\quad +x_{124}
\end{align*}
\begin{align*}
r_{107}&=-1+x_{120}^{2}+x_{123}^{2}-x_{124}^{2}
\end{align*}
\begin{align*}
r_{108}&=-17-8x_{4}-x_{4}^{2}+2x_{4}x_{120}+2x_{4}x_{121}-2x_{4}x_{171}+8x_{120}\\
&\quad +2x_{120}x_{171}+8x_{121}+2x_{121}x_{171}-x_{125}x_{126}-8x_{171}-x_{171}^{2}
\end{align*}
\begin{align*}
r_{109}&=-x_{21}-x_{21}x_{214}-x_{21}x_{215}+x_{126}
\end{align*}
\begin{align*}
r_{110}&=-4x_{128}-4x_{128}x_{133}+x_{134}
\end{align*}
\begin{align*}
r_{111}&=-x_{120}x_{128}-x_{121}x_{128}+x_{135}
\end{align*}
\begin{align*}
r_{112}&=-1+x_{127}^{2}+x_{128}^{2}-x_{135}^{2}
\end{align*}
\begin{align*}
r_{113}&=-x_{120}x_{130}-x_{121}x_{130}+x_{136}
\end{align*}
\begin{align*}
r_{114}&=-1+x_{129}^{2}+x_{130}^{2}-x_{136}^{2}
\end{align*}
\begin{align*}
r_{115}&=-x_{132}-x_{132}x_{134}+x_{137}
\end{align*}
\begin{align*}
r_{116}&=-1+x_{131}^{2}+x_{132}^{2}-x_{137}^{2}
\end{align*}
\begin{align*}
r_{117}&=-x_{128}^{2}+x_{139}
\end{align*}
\begin{align*}
r_{118}&=x_{130}-x_{138}x_{139}-x_{139}
\end{align*}
\begin{align*}
r_{119}&=1-x_{120}-x_{121}-x_{129}x_{140}+x_{134}
\end{align*}
\begin{align*}
r_{120}&=-x_{127}-x_{129}x_{141}+x_{131}
\end{align*}
\begin{align*}
r_{121}&=-25-x_{12}-13x_{13}-4x_{23}-2x_{26}-x_{30}-x_{32}-x_{34}-x_{36}-2x_{40}\\
&\quad -4x_{42}-2x_{43}-x_{46}-2x_{68}-4x_{128}x_{142}+x_{132}
\end{align*}
\begin{align*}
r_{122}&=-x_{143}x_{144}+x_{145}
\end{align*}
\begin{align*}
r_{123}&=-1-2x_{120}x_{145}-2x_{121}x_{145}+x_{143}^{2}+x_{144}^{2}
\end{align*}
\begin{align*}
r_{124}&=-2-2x_{3}-x_{3}^{2}+2x_{3}x_{120}+2x_{3}x_{121}+2x_{120}+2x_{121}-x_{146}
\end{align*}
\begin{align*}
r_{125}&=x_{3}x_{143}-x_{5}+x_{143}-x_{144}-x_{146}x_{147}
\end{align*}
\begin{align*}
r_{126}&=-x_{148}x_{149}+x_{150}
\end{align*}
\begin{align*}
r_{127}&=-1-2x_{120}x_{150}-2x_{121}x_{150}+x_{148}^{2}+x_{149}^{2}
\end{align*}
\begin{align*}
r_{128}&=-5-4x_{4}-x_{4}^{2}+2x_{4}x_{120}+2x_{4}x_{121}+4x_{120}+4x_{121}-x_{151}
\end{align*}
\begin{align*}
r_{129}&=x_{4}x_{148}-x_{21}+2x_{148}-x_{149}-x_{151}x_{152}
\end{align*}
\begin{align*}
r_{130}&=-1-x_{21}^{2}+2x_{21}x_{120}+2x_{21}x_{121}-x_{153}
\end{align*}
\begin{align*}
r_{131}&=-3-x_{4}+x_{21}x_{148}-x_{149}-x_{153}x_{154}-x_{171}
\end{align*}
\begin{align*}
r_{132}&=-10-6x_{4}-x_{4}^{2}+2x_{4}x_{120}+2x_{4}x_{121}-2x_{4}x_{171}+6x_{120}\\
&\quad +2x_{120}x_{171}+6x_{121}+2x_{121}x_{171}-x_{155}-6x_{171}-x_{171}^{2}
\end{align*}
\begin{align*}
r_{133}&=-1+x_{4}x_{143}+3x_{143}+x_{143}x_{171}-x_{144}-x_{155}x_{156}-x_{169}\\
&\quad -x_{170}
\end{align*}
\begin{align*}
r_{134}&=-x_{157}x_{158}+x_{159}
\end{align*}
\begin{align*}
r_{135}&=-1-2x_{120}x_{159}-2x_{121}x_{159}+x_{157}^{2}+x_{158}^{2}
\end{align*}
\begin{align*}
r_{136}&=-1+x_{4}x_{157}-x_{155}x_{160}+3x_{157}+x_{157}x_{171}-x_{158}-x_{201}\\
&\quad -x_{202}
\end{align*}
\begin{align*}
r_{137}&=-x_{161}x_{162}+x_{163}
\end{align*}
\begin{align*}
r_{138}&=-1-2x_{120}x_{163}-2x_{121}x_{163}+x_{161}^{2}+x_{162}^{2}
\end{align*}
\begin{align*}
r_{139}&=-1+x_{4}x_{161}-x_{155}x_{164}+3x_{161}+x_{161}x_{171}-x_{162}-x_{214}\\
&\quad -x_{215}
\end{align*}
\begin{align*}
r_{140}&=-x_{165}x_{166}+x_{167}
\end{align*}
\begin{align*}
r_{141}&=-1-2x_{120}x_{167}-2x_{121}x_{167}+x_{165}^{2}+x_{166}^{2}
\end{align*}
\begin{align*}
r_{142}&=-2-x_{4}+x_{4}x_{165}-x_{4}x_{169}-x_{4}x_{170}+4x_{165}+x_{165}x_{171}\\
&\quad -x_{166}-x_{168}-x_{168}x_{176}-x_{168}x_{177}-3x_{169}-2x_{170}
\end{align*}
\begin{align*}
r_{143}&=-x_{172}x_{173}+x_{174}
\end{align*}
\begin{align*}
r_{144}&=-1-2x_{120}x_{174}-2x_{121}x_{174}+x_{172}^{2}+x_{173}^{2}
\end{align*}
\begin{align*}
r_{145}&=-x_{21}-x_{21}x_{176}-x_{21}x_{177}+x_{178}
\end{align*}
\begin{align*}
r_{146}&=x_{4}x_{172}+x_{171}x_{172}+4x_{172}-x_{173}-x_{175}x_{178}-x_{176}
\end{align*}
\begin{align*}
r_{147}&=-x_{179}x_{180}+x_{181}
\end{align*}
\begin{align*}
r_{148}&=-1-2x_{120}x_{181}-2x_{121}x_{181}+x_{179}^{2}+x_{180}^{2}
\end{align*}
\begin{align*}
r_{149}&=x_{4}x_{179}+x_{171}x_{179}-x_{178}x_{182}+4x_{179}-x_{180}-x_{183}
\end{align*}
\begin{align*}
r_{150}&=x_{176}+x_{177}-x_{183}-x_{184}
\end{align*}
\begin{align*}
r_{151}&=-x_{185}x_{186}+x_{187}
\end{align*}
\begin{align*}
r_{152}&=-1-2x_{120}x_{187}-2x_{121}x_{187}+x_{185}^{2}+x_{186}^{2}
\end{align*}
\begin{align*}
r_{153}&=x_{4}x_{185}+x_{171}x_{185}-x_{178}x_{188}+4x_{185}-x_{186}-x_{189}
\end{align*}
\begin{align*}
r_{154}&=x_{176}+x_{177}-x_{189}-x_{190}
\end{align*}
\begin{align*}
r_{155}&=-x_{191}x_{192}+x_{193}
\end{align*}
\begin{align*}
r_{156}&=-1-2x_{120}x_{193}-2x_{121}x_{193}+x_{191}^{2}+x_{192}^{2}
\end{align*}
\begin{align*}
r_{157}&=x_{4}x_{191}+x_{171}x_{191}-x_{178}x_{194}+4x_{191}-x_{192}-x_{195}
\end{align*}
\begin{align*}
r_{158}&=x_{176}+x_{177}-x_{195}-x_{196}
\end{align*}
\begin{align*}
r_{159}&=-x_{197}x_{198}+x_{199}
\end{align*}
\begin{align*}
r_{160}&=-1-2x_{120}x_{199}-2x_{121}x_{199}+x_{197}^{2}+x_{198}^{2}
\end{align*}
\begin{align*}
r_{161}&=-x_{21}-x_{21}x_{201}-x_{21}x_{202}+x_{203}
\end{align*}
\begin{align*}
r_{162}&=x_{4}x_{197}+x_{171}x_{197}+4x_{197}-x_{198}-x_{200}x_{203}-x_{201}
\end{align*}
\begin{align*}
r_{163}&=-x_{204}x_{205}+x_{206}
\end{align*}
\begin{align*}
r_{164}&=-1-2x_{120}x_{206}-2x_{121}x_{206}+x_{204}^{2}+x_{205}^{2}
\end{align*}
\begin{align*}
r_{165}&=x_{4}x_{204}+x_{171}x_{204}-x_{203}x_{207}+4x_{204}-x_{205}-x_{208}
\end{align*}
\begin{align*}
r_{166}&=x_{201}+x_{202}-x_{208}-x_{209}
\end{align*}
\begin{align*}
r_{167}&=-x_{210}x_{211}+x_{212}
\end{align*}
\begin{align*}
r_{168}&=-1-2x_{120}x_{212}-2x_{121}x_{212}+x_{210}^{2}+x_{211}^{2}
\end{align*}
\begin{align*}
r_{169}&=x_{4}x_{210}-x_{126}x_{213}+x_{171}x_{210}+4x_{210}-x_{211}-x_{214}
\end{align*}
\begin{align*}
r_{170}&=-x_{216}x_{217}+x_{218}
\end{align*}
\begin{align*}
r_{171}&=-1-2x_{120}x_{218}-2x_{121}x_{218}+x_{216}^{2}+x_{217}^{2}
\end{align*}
\begin{align*}
r_{172}&=x_{4}x_{216}-x_{126}x_{219}+x_{171}x_{216}+4x_{216}-x_{217}-x_{220}
\end{align*}
\begin{align*}
r_{173}&=x_{214}+x_{215}-x_{220}-x_{221}
\end{align*}
\begin{align*}
r_{174}&=-x_{222}x_{223}+x_{224}
\end{align*}
\begin{align*}
r_{175}&=-1-2x_{120}x_{224}-2x_{121}x_{224}+x_{222}^{2}+x_{223}^{2}
\end{align*}
\begin{align*}
r_{176}&=-x_{21}-x_{21}x_{169}-x_{21}x_{170}+x_{228}
\end{align*}
\begin{align*}
r_{177}&=x_{4}x_{222}+x_{171}x_{222}+4x_{222}-x_{223}-x_{225}x_{228}-x_{226}
\end{align*}
\begin{align*}
r_{178}&=x_{169}+x_{170}-x_{226}-x_{227}
\end{align*}
\begin{align*}
r_{179}&=x_{4}x_{128}+x_{127}+4x_{128}+x_{128}x_{171}-x_{135}-x_{178}x_{229}-x_{230}
\end{align*}
\begin{align*}
r_{180}&=x_{176}+x_{177}-x_{230}-x_{231}
\end{align*}
\begin{align*}
r_{181}&=-26-x_{12}-13x_{13}-4x_{23}-2x_{26}-x_{30}-x_{32}-x_{34}-x_{36}-2x_{40}\\
&\quad -4x_{42}-2x_{43}-x_{46}-2x_{68}+x_{128}-x_{143}-x_{148}-x_{157}-x_{161}\\
&\quad -x_{165}-x_{172}-x_{179}-x_{185}-x_{191}-x_{197}-x_{204}-x_{210}\\
&\quad -x_{216}-x_{222}-x_{232}
\end{align*}
\begin{align*}
r_{182}&=1970+216x_{0}+6x_{0}^{2}+8x_{0}x_{1}+16x_{0}x_{2}+2x_{0}x_{3}+8x_{0}x_{9}\\
&\quad -2x_{0}x_{143}-4x_{0}x_{148}-2x_{0}x_{165}+154x_{1}+5x_{1}^{2}+16x_{1}x_{2}+8x_{1}x_{9}\\
&\quad -4x_{1}x_{148}-2x_{1}x_{161}+308x_{2}+17x_{2}^{2}+18x_{2}x_{9}-8x_{2}x_{148}\\
&\quad -2x_{2}x_{157}+32x_{3}+x_{3}^{2}-2x_{3}x_{165}+156x_{9}+5x_{9}^{2}-4x_{9}x_{148}\\
&\quad -2x_{9}x_{157}+2x_{12}+x_{12}^{2}-2x_{12}x_{172}+2x_{13}+x_{13}^{2}-2x_{13}x_{179}\\
&\quad +2x_{14}+x_{14}^{2}-2x_{14}x_{185}+2x_{15}+x_{15}^{2}-2x_{15}x_{191}+2x_{16}+x_{16}^{2}\\
&\quad -2x_{16}x_{197}+2x_{17}+x_{17}^{2}-2x_{17}x_{204}+2x_{18}+x_{18}^{2}-2x_{18}x_{210}\\
&\quad +2x_{19}+x_{19}^{2}-2x_{19}x_{216}+2x_{20}+x_{20}^{2}-2x_{20}x_{222}-x_{120}x_{233}\\
&\quad -x_{121}x_{233}-32x_{143}+x_{143}^{2}-76x_{148}+x_{148}^{2}-4x_{157}+x_{157}^{2}\\
&\quad -2x_{161}+x_{161}^{2}-32x_{165}+x_{165}^{2}-2x_{172}+x_{172}^{2}-2x_{179}+x_{179}^{2}\\
&\quad -2x_{185}+x_{185}^{2}-2x_{191}+x_{191}^{2}-2x_{197}+x_{197}^{2}-2x_{204}+x_{204}^{2}\\
&\quad -2x_{210}+x_{210}^{2}-2x_{216}+x_{216}^{2}-2x_{222}+x_{222}^{2}+x_{233}
\end{align*}
\endgroup


\end{document}
